\ArithmeticHead{384}{L}{ANALYSIS.\quad NACHLASS.}
\begingroup
\makeatletter\tagsleft@true\let\veqno\@@leqno\makeatother
\begin{center}[7.]\par\medskip
\rule{.36\linewidth}{.4pt}\par\smallskip
{\small[Eintragungen im \textsc{Leiste}.]}\par\smallskip
\rule{.36\linewidth}{.4pt}
\end{center}
\noindent{\small[S. 58]}
\[
\begin{aligned}
1-x+\frac{xx}{4}-\frac{x^3}{36}+\text{etc.}&=\varphi\\[.5em]
x-\frac{x^2}{2}+\frac{x^3}{12}-\text{etc.}&=\psi\\[.5em]
\varphi'=-\psi:x,\qquad\psi'&=\varphi.
\end{aligned}
\]
\noindent{\small[S. 65]}
\[
z=1-x+\frac{x^2}{4}-\frac{x^3}{4{.}9}+{.}\,{.}\,{.}\ [=\varphi]
\]
\[
x\frac{dz}{dx}=-x+\frac{x^2}{2}-\frac{x^3}{4{.}3}+\frac{x^4}{4{.}9{.}4}\text{etc.}\ [=-\psi]
\]
\begin{equation}
x\frac{d^2z}{dx^2}+\frac{dz}{dx}+z=0.\tag*{[12]}
\end{equation}
\begin{equation}
\begin{aligned}[t]
z=yx^m,\qquad\frac{dz}{dx}&=mx^{m-1}y+x^m\frac{dy}{dx}\\[.7em]
\frac{d^2y}{dx^2}+\frac{2m+1}{x}\frac{dy}{dx}&+\left(\frac{mm}{xx}+\frac1x\right)y\ [=0]
\end{aligned}\tag*{[13]}
\end{equation}
\clearpage\SourceLeaf{392}\ArithmeticHead{385}{R}{SUMMATION EINER DIVERGIRENDEN REIHE.}
\begin{equation}
m=-\tfrac12\tag*{[14]}
\end{equation}
\begin{equation}
\frac{d^2y}{dx^2}+\left(\frac14\,\frac1{x^2}+\frac1x\right)y\ [=0].\tag*{[15]}
\end{equation}
\ClosingRule
\begin{center}[8.]\par\medskip
\rule{.36\linewidth}{.4pt}\par\smallskip
{\small[Aus dem \textit{Tagebuch}, Nr. 82.]}\par\smallskip
\rule{.36\linewidth}{.4pt}
\end{center}
\noindent\hspace*{\parindent}Seriei
\begin{equation}
x-\frac12x^2+\frac1{12}x^3-\frac1{144}x^4\,{.}\,{.}\,{.}\ [=\psi]\tag*{[16]}
\end{equation}
\noindent summam consideravimus invenimusque eam \(=0\), si
\begin{equation}
2\surd x+\frac3{16}\,\frac1{\surd x}-\frac{21}{1024}\,\frac1{\surd x^3}\,{.}\,{.}\,{.}=\left(k+\frac14\right)\pi\text{[\textsuperscript{*})]}\tag*{[17]}
\end{equation}
\begin{flushright}Brunsv[igae,] Oct. 16. [1797]\end{flushright}
\par\medskip\noindent\rule{.24\linewidth}{.4pt}\par\smallskip
{\small\noindent\hspace*{\parindent}[\textsuperscript{*}) In der Handschrift steht links vom Gleichheitszeichen im dritten Gliede \(\surd{.}\,3x\) statt \(\surd x^3\); vergl. in den \textls[120]{Bemerkungen} S. 389 die Bestätigung der obigen Formel]\par}
\par\endgroup\ArithmeticHead{385}{R}{SUMMATION EINER DIVERGIRENDEN REIHE.}
